\scanpage{65}

On peut prendre pour $\zeta$ n'importe quelle racine de $1$, à l'exclusion de $\pm1$ (racines d'ordre $1$ et $2$ \dots).

Ainsi on a trouvé

\uline{Proposition} a) Tout \deleted{\unsure{épinglé|choix}} polygône régulier épinglé à $n$ cotés, $3\leq n<+\infty$ (donc $n$ fini), sur un corps alg. clos $k$ de car. $p$ premier à $n$, est isomorphe à un \inserted{des} polygônes épinglés $(\Pi_{\zeta,k},\deleted{\unsure{s}}\,\unsure{r_{n,k}|r_{\zeta,k}})$ \inserted{dans $k^2$} construits précédemment, \inserted{où $\zeta$ est racine primitive $n$ ième de $1$.} Pour \unsure{$\Pi,r$|$(\Pi,r)$|$\Pi,n$|$\Pi$} fixé, $\zeta$ est défini modulo passage à $\zeta^{-1}$, \deleted{\unsure{on}}

b) On trouve ainsi une correspondance $1$-$1$ entre classes d'isom. de $n$-polygônes réguliers sur $k$ (ou classes d'isom de $n$-polygônes réguliers épinglés), entre \deleted{couples} \inserted{paires} $\{\zeta,\zeta^{-1}\}$ de paires de racines \unsure{primes|primièmes|primitives} $n$\textsuperscript{ièmes} de $1$ inverses l'une de l'autre, et entre éléments $\alpha\in k$ tels que les \inserted{$2$} solutions de l'équation

\margin{terminologie: \uline{invariant} \uline{numérique} des polygônes réguliers \dots}
% slanted

\[\zeta^2\deleted{\unsure{+}}-\alpha\zeta+1=0\]

soient des racines primitives $n$\textsuperscript{ièmes} de $1$.

c) Dans $\Pi_{\zeta,k}$, l'origine de $k^2$ est \unsure{car.|aussi} comme le barycentre de $S$, ou encore l'unique pt fixe sous $G$\inserted{$=\operatorname{Aut}(\Pi_{\unsure{n|m|\zeta},k})$}, ou l'unique pt fixe sous $G^+$. $\{\zeta,\zeta^{-1}\}$ est l'ens des valeurs

% archive sheet number: 64
